\section{Experiments}\label{sec:Experiments}
In this section, we empirically study the behavior of $1$-GNNs, $2$-GNNs, and $3$-GNNs on the problem of distinguishing graphs that contain a $k$-clique from graphs that are $(k\!-\!1)$-colorable. 
% We evaluate both training and test performance, using the former to assess how well the models fit the task and the latter to measure their ability to generalize.

We evaluate the three GNN architectures under two different experimental settings. In Experiment~1, we train and test the models on \emph{planted random graphs}. These are random $N$-vertex graphs generated either by planting a $k$-clique, or by imposing a $(k\!-\!1)$-partition in which each color class forms an independent set and edges are added independently across classes. In Experiment~2, we modify this data-generation procedure so that the two graph classes have matching average degree (this adjustment is described in detail in Appendix ~\ref{app:datGeneration}). 
% In Experiment~3, we evaluate the models on the hard graph family $\Fam_k$ introduced in Section~\ref{sec:HardFamilies}.


\subsection{Dataset generation}\label{subsec:DataGeneration}

We consider two different graph-generation procedures corresponding to the experimental settings described above.

\textbf{Planted random graphs.} 
In the first experimental setting, graphs are generated uniformly at random from two classes: graphs on $N$ vertices that contain a $k$-clique, and graphs on $N$ vertices that are $(k\!-\!1)$-colorable. Both classes are generated using Erd\H{o}s--R\'enyi-style random constructions with edge probability $p=0.5$. In particular, $(k\!-\!1)$-colorable graphs are generated by randomly partitioning the vertex set into $(k\!-\!1)$ parts and independently adding edges between vertices in different parts with probability $0.5$, while $k$-clique graphs are generated by first planting a $k$-clique and then adding all remaining edges independently at random with probability $0.5$. As node features, we use the degree of each vertex.



% \begin{observation}[Expected average degree gap in Experiment~1]
% \label{obs:avg-degree-gap}
% Fix $N \in \N$ and $k \ge 3$, and let $p=0.5$.
% \begin{enumerate}
%     \item (\emph{$(k\!-\!1)$-colorable graphs.})
%     Generate a random $(k\!-\!1)$-colorable graph by partitioning $V$ uniformly at random into $(k\!-\!1)$ parts and then adding each edge between different parts independently with probability $p$.
%     Then the expected average degree satisfies
%     \[
%     \E[d_{\avg}] \;=\; p\,(N-1)\frac{k-2}{k-1}
%     \;=\; \frac{1}{2}(N-1)\frac{k-2}{k-1}.
%     \]
%     \item (\emph{Graphs with a planted $k$-clique.})
%     Generate a random graph with a planted $k$-clique by first selecting a $k$-subset $S \subseteq V$ and adding all edges inside $S$, and then adding every remaining edge independently with probability $p$.
%     Then the expected average degree satisfies
%     \[
%     \E[d_{\avg}]
%     \;=\;
%     \frac{k(N-1) + (N-k)(N-k-1)}{N}.
%     \]
% \end{enumerate}
% In particular, for $N>k$ one has $\E[d_{\avg}]$ larger in the planted-$k$-clique model than in the $(k\!-\!1)$-colorable model, yielding an explicit average-degree gap.
% \end{observation}






\textbf{Matching average degree.}
To eliminate trivial separation by degree, the second experimental setting modifies the data-generation process so that the two graph classes have matching expected average degree. This is achieved by adjusting the edge probability in the random construction of the $k$-clique graphs so that their expected average degree coincides with that of the $(k\!-\!1)$-colorable graphs. We give the details of this adjustment in appendix \ref{app:datGeneration}.

This setting controls for degree-based heuristics and provides a more challenging benchmark for assessing whether GNNs can learn more subtle structural features.

% \textbf{Hard graph family.}
% In the third experimental setting, we evaluate the models on the explicit hard graph family $\Fam_k$ introduced in Section~\ref{sec:HardFamilies}.


\subsection{Models and configurations}\label{sec:ModelsandConfig}
We evaluate three classes of GNN architectures following the higher-order framework of \citet{MorrisEtAl2019}: standard 1-GNNs (GCN-based message-passing networks with 1-WL expressivity), hierarchical 1-2-GNNs (2-GNNs initialized with 1-GNN features), and hierarchical 1-3-GNNs (3-GNNs initialized with 1-GNN features). All models use degree-based node features, global mean pooling, and an MLP classifier, and are trained with AdamW and early stopping. To ensure a fair and comprehensive comparison across experiments, we cast a wide architectural net and perform grid searches over hidden dimensions, pooling choices, and learning rates for each model class. The higher-order GNNs use the \textsc{TwoMalkin} and \textsc{ConnectedThreeMalkin} constructions to lift representations to the 2-local and 3-local graphs, respectively, and concatenate pooled embeddings across levels. Full architectural specifications, hyperparameter grids, and training details are provided in Appendix~\ref{app:Experiments}.



\subsection{Results and Discussion}\label{subsec:ResultsandDiscussion}
In Experiment~1, 
% we evaluated the performance of 1-GNNs and 2-GNNs on planted random graphs, where graphs from the positive class contain a planted $k$-clique and graphs from the negative class are $(k-1)$-colorable, both generated with Erdős--Rényi edge probability $p=0.5$. Due to the additional edges introduced by the planted clique, the two classes exhibit noticeably different degree distributions. As a result, 
both 1-GNNs and 2-GNNs achieved perfect test accuracy (100\%) across all tested clique sizes $k \in \{3,4,5,6,7\}$ and graph sizes $N \in \{100,\dots,800\}$ (with 2-GNNs evaluated at $N=100$), demonstrating that the task is trivial for standard GNNs when degree signals are present (see table\ref{tab:exp1_results}).

In Experiment~2, we removed the degree shortcut by matching the expected average degree between the planted $k$-clique and $(k\!-\!1)$-colorable classes. Under this harder setting, 1-GNN accuracy dropped sharply with increasing $N$ and $k$, often approaching chance level, though occasional high scores indicate that the task remains learnable but difficult to optimize. These results show that once degree information is neutralized, standard GNNs struggle to detect the planted structure reliably.

The main results of experiment 1 and 2 are summarized in Table \ref{tab:exp1_and_exp2}.


% \begin{table}[ht]
% \centering
% \caption{Test accuracy (\%) on the planted random graphs setting (Experiment 1) for different clique sizes $k$.}
% \label{tab:exp1_results}
% \begin{tabular}{lcc}
% \toprule
% $k$ & 1-GNN & 2-GNN \\
% \midrule
% 3   & 100 & 100 \\
% 4   & 100 & 100 \\
% 5   & 100 & 100 \\
% 6   & 100 & 100 \\
% 7   & 100 & 100 \\
% \bottomrule
% \end{tabular}
% \end{table}

% \begin{table}[ht]
% \centering
% \caption{Exeperiment 2, degree-matched graphs ($N=100$)}
% \label{tab:wl_comparison_n100}
% \begin{tabular}{lccc}
% \toprule
% $k$ & 1-GNN  & 2-GNN & 3-GNN  \\
% \midrule
% 3   & 68\%         & \textbf{95\%} & 93\%          \\
% 4   & \textbf{85\%}& 64\%          & 57\%          \\
% 5   & \textbf{75\%}& 66\%          & 59\%          \\
% 6   & \textbf{62\%}& 56\%          & 53\%          \\
% 7   & \textbf{63\%}& 55\%          & —             \\
% \bottomrule
% \end{tabular}
% \end{table}

\begin{table}[ht]
\centering
\caption{Results on planted graphs and degree-matched setting (at $N=100$) .}
\label{tab:exp1_and_exp2}

\begin{subtable}{0.48\textwidth}
\centering
\caption{Planted Random Graphs}
\label{tab:exp1_results}
\begin{tabular}{lcc}
\toprule
$k$ & 1-GNN & 2-GNN \\
\midrule
3   & 100 & 100 \\
4   & 100 & 100 \\
5   & 100 & 100 \\
6   & 100 & 100 \\
7   & 100 & 100 \\
\bottomrule
\end{tabular}
\end{subtable}
\hfill
\begin{subtable}{0.48\textwidth}
\centering
\caption{Degree matched random graphs}
\label{tab:wl_comparison_n100}
\begin{tabular}{lrrr}
\toprule
$k$ & 1-GNN  & 2-GNN  & 3-GNN \\
\midrule
3   & 68  & \textbf{95} & 93  \\
4   & \textbf{85}& 64          & 57  \\
5   & \textbf{75}& 66          & 59  \\
6   & \textbf{62}& 56          & 53  \\
7   & \textbf{63}& 55          & —   \\
\bottomrule
\end{tabular}
\end{subtable}

\end{table}

% In Experiment~3, we evaluated GNNs on the hard graph family $(G^3_s, H^3_s)$, which share the same size and degree sequence and are provably indistinguishable by both 1-WL and 2-WL. As predicted, 1-GNNs and 2-GNNs achieved only random-guessing accuracy (about 50\%), fully matching the theoretical impossibility results. In contrast, a 3-GNN was able to learn the distinction and reached up to 98\% accuracy, indicating that higher-order message passing is necessary to separate this graph family.

% Results of experiment 3 are summarized in Table \ref{tab:exp3_hard_family}.


% \begin{table}[ht]
% \centering
% \caption{Performance on the explicit hard family $\mathcal{F}_3$ (Experiment 3).}
% \label{tab:exp3_hard_family}
% \begin{tabular}{lccc}
% \toprule
% Model & WL Power & Best Val. Accuracy & Theory Prediction \\
% \midrule
% \textbf{1-GNN} & 1-WL & \textbf{52\%} ($\approx$ random) & Cannot distinguish (Theorem 3.2) \\
% \textbf{2-GNN} & 2-WL & \textbf{54\%} ($\approx$ random) & Cannot distinguish (Theorem 3.3) \\
% \textbf{3-GNN} & 3-WL & \textbf{98\%}                    & No hardness theorem for 3-WL on this family \\
% \bottomrule
% \end{tabular}
% \end{table}
